% !PRINTED 639
\weberpage{639}
9. \gesperrt{Ein aus einer Minimalbasis $\omega_1, \omega_2, \ldots, \omega_n$ abgeleitetes
 System ganzer Funktionen}

\weq{7}{\omega_\nu' = x_{\nu,1} \omega_1 + x_{\nu,2} \omega_2 + \cdots + x_{\nu,n} \omega_n,}
\[
\nu = 1, 2, \ldots, n
\]

\gesperrt{ist dann und nur dann eine Minimalbasis, wenn die
Determinante}

\weq{8}{X = \Sigma \pm x_{1,1} x_{2,2} \ldots x_{n,n}}

\gesperrt{eine von Null verschiedene Konstante ist.}

Denn hat diese Determinante, die eine rationale Funktion
von $z$ ist, einen Linearfaktor $z - c$, so kann man die Konstanten
$c_1, c_2, \ldots, c_n$ so bestimmen, daß die $n$ ganzen rationalen Funktionen


\[
c_1 x_{1,\nu} + c_2 x_{2,\nu} + \cdots + c_n x_{n,\nu}
\]

durch $z - c$ teilbar werden, ohne daß alle $c_i$ verschwinden.
Dann ist aber

\[
\frac{c_1 \omega_1' + c_2 \omega_2' + \cdots + c_n \omega_n'}{z - c}
\]

eine ganze Funktion, und $\omega_1', \omega_2', \ldots, \omega_n'$ keine Minimalbasis.

Für die Diskriminante erhält man aus (7) nach § 173, (6)

\weq{9}{\varDelta(\omega_1', \omega_2', \ldots, \omega_n') = X^2 \varDelta(\omega_1, \omega_2, \ldots, \omega_n).}

Daraus folgt:

10. \gesperrt{Die Diskriminante einer Minimalbasis ist, von
einem konstanten Faktor abgesehen, von der besonderen
 Wahl der Basis unabhängig.}

Die Diskriminante einer Minimalbasis hat unter allen Diskriminanten
 aus ganzen Funktionen gebildeter Basen den \gesperrt{niedrigsten
 Grad} (daher der Name Minimalbasis). Es ist eine durch
den Körper selbst, abgesehen von einem konstanten Faktor,
eindeutig bestimmte ganze rationale Funktion von $z$. Sie wird
daher auch \gesperrt{die Diskriminante des Körpers $\Omega$} genannt und
mit $\varDelta(\Omega)$ bezeichnet.
